%
% display-server.tex
%
% Z Specification for the Lux Display Server
% Formal model of DisplayServer._on_frame() and FrameReader
%

\documentclass[a4paper,10pt,fleqn]{article}
\usepackage[margin=1in]{geometry}
\usepackage{fuzz}
\usepackage[colorlinks=true,linkcolor=blue,citecolor=blue,urlcolor=blue]{hyperref}

\begin{document}

\title{Lux Display Server: A Z Specification}
\author{Formal Model of Frame Lifecycle, Client Management, and Scene State}
\date{March 2026}
\hypersetup{
  pdftitle={Lux Display Server: A Z Specification},
  pdfauthor={Punt Labs},
  pdfsubject={Z Specification},
  pdfcreator={fuzz/probcli}
}
\maketitle

\tableofcontents
\newpage

% ===================================================================
\section{Introduction}
% ===================================================================

The Lux display server is an ImGui render loop that listens on a Unix
domain socket and renders scenes.  Each frame, it accepts new
connections, polls clients for messages, renders the current scene,
and flushes interaction events.

This specification models the stateful core:
\begin{itemize}
  \item Client connection tracking and FrameReader association
  \item Scene replacement, incremental patching, and clearing
  \item Message dispatch (scene, update, clear, ping)
  \item Interaction event generation and broadcast
  \item The FrameReader streaming buffer
\end{itemize}

We abstract away ImGui rendering (it is stateless from the
protocol's perspective) and focus on the state that the protocol
cares about.  Element ordering and event queue ordering are
abstracted to sets for ProB animatability --- $\seq$ types
cause unbounded enumeration during model checking.

% ===================================================================
\section{Basic Types}
% ===================================================================

\begin{zed}
  [FD, ELEMID, SCENEID, RAWBYTES, ACTION]
\end{zed}

$FD$ models socket file descriptors.  $ELEMID$ and $SCENEID$ are
opaque identifiers for elements and scenes.  $RAWBYTES$ represents
a chunk of bytes received from a socket.  $ACTION$ identifies
the action string on a button click.

% ===================================================================
\section{Free Types}
% ===================================================================

% Boolean type (fuzz doesn't have built-in boolean)
\begin{zed}
  ZBOOL ::= ztrue | zfalse
\end{zed}

Element kinds supported in Phase~1:
\begin{zed}
  ElementKind ::= ekText | ekButton | ekSeparator | ekImage
\end{zed}

Client message types:
\begin{zed}
  ClientMsgType ::= cmScene | cmUpdate | cmClear | cmPing
\end{zed}

Display message types (sent back to clients):
\begin{zed}
  DisplayMsgType ::= dmReady | dmAck | dmPong | dmInteraction
\end{zed}

% ===================================================================
\section{Global Constants}
% ===================================================================

% Animation-friendly bounds.  Production values (64, 1000, 256, 16M)
% are too large for ProB's state space --- reduce for model checking.
\begin{axdef}
  maxClients : \nat \\
  maxElements : \nat \\
  maxEvents : \nat \\
  maxBufSize : \nat
  \where
  maxClients = 3 \\
  maxElements = 3 \\
  maxEvents = 3 \\
  maxBufSize = 4
\end{axdef}

% ===================================================================
\section{Display Server State}
% ===================================================================

The server tracks connected clients, the current scene (if any),
a set of pending interaction events, and the FrameReader buffer
state.  All fields are flattened into a single state schema for
ProB animatability.

\begin{schema}{DisplayServer}
  clients : \power FD \\
  readers : \power FD \\
  hasScene : ZBOOL \\
  sceneId : SCENEID \\
  elemIds : \power ELEMID \\
  elemKinds : ELEMID \pfun ElementKind \\
  eventQueue : \power ELEMID \\
  listening : ZBOOL \\
  bufSize : \nat \\
  pendingMsgs : \nat
  \where
  readers = clients \\
  \# clients \leq maxClients \\
  \# elemIds \leq maxElements \\
  \# eventQueue \leq maxEvents \\
  bufSize \leq maxBufSize \\
  elemIds \subseteq \dom elemKinds \\
  hasScene = ztrue \implies
\quad~ eventQueue \subseteq elemIds
\end{schema}

The key invariants are:
\begin{enumerate}
  \item Every client has a FrameReader ($readers = clients$).
  \item Client count, element count, event queue, and buffer are bounded.
  \item Element kinds are defined for all elements in the scene.
  \item Queued events reference elements in the current scene.
\end{enumerate}

% ===================================================================
\section{Initialization}
% ===================================================================

\begin{schema}{Init}
  DisplayServer'
  \where
  clients' = \emptyset \\
  readers' = \emptyset \\
  hasScene' = zfalse \\
  elemIds' = \emptyset \\
  elemKinds' = \emptyset \\
  eventQueue' = \emptyset \\
  listening' = ztrue \\
  bufSize' = 0 \\
  pendingMsgs' = 0
\end{schema}

% ===================================================================
\section{Operations}
% ===================================================================

% -------------------------------------------------------------------
\subsection{Accept Connection}
% -------------------------------------------------------------------

A new client connects.  We allocate a fresh FrameReader for it.

\begin{schema}{AcceptConnection}
  \Delta DisplayServer \\
  newClient? : FD
  \where
  listening = ztrue \\
  newClient? \notin clients \\
  \# clients < maxClients \\
  clients' = clients \cup \{ newClient? \} \\
  readers' = readers \cup \{ newClient? \} \\
  hasScene' = hasScene \\
  sceneId' = sceneId \\
  elemIds' = elemIds \\
  elemKinds' = elemKinds \\
  eventQueue' = eventQueue \\
  listening' = listening \\
  bufSize' = bufSize \\
  pendingMsgs' = pendingMsgs
\end{schema}

% -------------------------------------------------------------------
\subsection{Disconnect Client}
% -------------------------------------------------------------------

A client disconnects (error or clean close).

\begin{schema}{DisconnectClient}
  \Delta DisplayServer \\
  deadClient? : FD
  \where
  deadClient? \in clients \\
  clients' = clients \setminus \{ deadClient? \} \\
  readers' = readers \setminus \{ deadClient? \} \\
  hasScene' = hasScene \\
  sceneId' = sceneId \\
  elemIds' = elemIds \\
  elemKinds' = elemKinds \\
  eventQueue' = eventQueue \\
  listening' = listening \\
  bufSize' = bufSize \\
  pendingMsgs' = pendingMsgs
\end{schema}

% -------------------------------------------------------------------
\subsection{Receive Scene}
% -------------------------------------------------------------------

A client sends a new scene, replacing any existing one.
The event queue is cleared since old element references are stale.

\begin{schema}{ReceiveScene}
  \Delta DisplayServer \\
  newSceneId? : SCENEID \\
  newElemIds? : \power ELEMID \\
  newElemKinds? : ELEMID \pfun ElementKind
  \where
  newElemIds? \subseteq \dom newElemKinds? \\
  \# newElemIds? \leq maxElements \\
  hasScene' = ztrue \\
  sceneId' = newSceneId? \\
  elemIds' = newElemIds? \\
  elemKinds' = newElemKinds? \\
  eventQueue' = \emptyset \\
  clients' = clients \\
  readers' = readers \\
  listening' = listening \\
  bufSize' = bufSize \\
  pendingMsgs' = pendingMsgs
\end{schema}

% -------------------------------------------------------------------
\subsection{Clear Scene}
% -------------------------------------------------------------------

A client sends a clear message, removing all content.

\begin{schema}{ClearScene}
  \Delta DisplayServer
  \where
  hasScene' = zfalse \\
  sceneId' = sceneId \\
  elemIds' = elemIds \\
  elemKinds' = elemKinds \\
  eventQueue' = \emptyset \\
  clients' = clients \\
  readers' = readers \\
  listening' = listening \\
  bufSize' = bufSize \\
  pendingMsgs' = pendingMsgs
\end{schema}

% -------------------------------------------------------------------
\subsection{Remove Element (Patch)}
% -------------------------------------------------------------------

An update message removes an element by ID from the current scene.

\begin{schema}{RemoveElement}
  \Delta DisplayServer \\
  targetId? : ELEMID
  \where
  hasScene = ztrue \\
  targetId? \in elemIds \\
  targetId? \notin eventQueue \\
  elemIds' = elemIds \setminus \{ targetId? \} \\
  elemKinds' = \{ targetId? \} \ndres elemKinds \\
  sceneId' = sceneId \\
  hasScene' = hasScene \\
  eventQueue' = eventQueue \\
  clients' = clients \\
  readers' = readers \\
  listening' = listening \\
  bufSize' = bufSize \\
  pendingMsgs' = pendingMsgs
\end{schema}

% -------------------------------------------------------------------
\subsection{Button Click (Interaction)}
% -------------------------------------------------------------------

The display detects a button click.  An interaction event is queued
for broadcast to all clients.

\begin{schema}{ButtonClick}
  \Delta DisplayServer \\
  buttonId? : ELEMID
  \where
  hasScene = ztrue \\
  buttonId? \in elemIds \\
  elemKinds~buttonId? = ekButton \\
  \# eventQueue < maxEvents \\
  eventQueue' = eventQueue \cup \{ buttonId? \} \\
  hasScene' = hasScene \\
  sceneId' = sceneId \\
  elemIds' = elemIds \\
  elemKinds' = elemKinds \\
  clients' = clients \\
  readers' = readers \\
  listening' = listening \\
  bufSize' = bufSize \\
  pendingMsgs' = pendingMsgs
\end{schema}

% -------------------------------------------------------------------
\subsection{Flush Events}
% -------------------------------------------------------------------

At the end of each frame, queued events are broadcast to all
clients and the queue is emptied.

\begin{schema}{FlushEvents}
  \Delta DisplayServer
  \where
  eventQueue' = \emptyset \\
  hasScene' = hasScene \\
  sceneId' = sceneId \\
  elemIds' = elemIds \\
  elemKinds' = elemKinds \\
  clients' = clients \\
  readers' = readers \\
  listening' = listening \\
  bufSize' = bufSize \\
  pendingMsgs' = pendingMsgs
\end{schema}

% -------------------------------------------------------------------
\subsection{Feed Bytes}
% -------------------------------------------------------------------

Feeding bytes into the FrameReader increases the buffer size.

\begin{schema}{FeedBytes}
  \Delta DisplayServer \\
  bytesIn? : \nat
  \where
  bytesIn? > 0 \\
  bytesIn? \leq maxBufSize \\
  bufSize + bytesIn? \leq maxBufSize \\
  bufSize' = bufSize + bytesIn? \\
  pendingMsgs' = pendingMsgs \\
  clients' = clients \\
  readers' = readers \\
  hasScene' = hasScene \\
  sceneId' = sceneId \\
  elemIds' = elemIds \\
  elemKinds' = elemKinds \\
  eventQueue' = eventQueue \\
  listening' = listening
\end{schema}

% -------------------------------------------------------------------
\subsection{Drain Messages}
% -------------------------------------------------------------------

Draining extracts all complete messages from the FrameReader,
reducing the buffer.

\begin{schema}{DrainMessages}
  \Delta DisplayServer \\
  drained! : \nat \\
  bytesConsumed? : \nat
  \where
  bytesConsumed? \leq maxBufSize \\
  bytesConsumed? \leq bufSize \\
  bufSize' = bufSize - bytesConsumed? \\
  pendingMsgs' = 0 \\
  drained! = pendingMsgs + bytesConsumed? \\
  clients' = clients \\
  readers' = readers \\
  hasScene' = hasScene \\
  sceneId' = sceneId \\
  elemIds' = elemIds \\
  elemKinds' = elemKinds \\
  eventQueue' = eventQueue \\
  listening' = listening
\end{schema}

% -------------------------------------------------------------------
\subsection{Shutdown}
% -------------------------------------------------------------------

The display server shuts down: all clients are disconnected,
the scene is cleared, events are flushed.

\begin{schema}{Shutdown}
  \Delta DisplayServer
  \where
  clients' = \emptyset \\
  readers' = \emptyset \\
  hasScene' = zfalse \\
  sceneId' = sceneId \\
  elemIds' = elemIds \\
  elemKinds' = elemKinds \\
  eventQueue' = \emptyset \\
  listening' = zfalse \\
  bufSize' = bufSize \\
  pendingMsgs' = pendingMsgs
\end{schema}

% ===================================================================
\section{System Invariants}
% ===================================================================

The key properties maintained across all operations:

\begin{enumerate}
  \item \textbf{Reader--client bijection}: Every connected client
    has a FrameReader, and every FrameReader belongs to
    a connected client ($readers = clients$).

  \item \textbf{Event validity}: Queued interaction events reference
    elements that exist in the current scene.

  \item \textbf{Bounded resources}: Client count, element count,
    event queue length, and buffer sizes are all bounded.

  \item \textbf{Scene atomicity}: A new scene completely replaces
    the old one and clears stale events.

  \item \textbf{Element kind coverage}: Every element ID in the
    scene has a corresponding kind in the kind map.
\end{enumerate}

\end{document}
